\scanpage{38}
(avec diagramme commutatifs).

\textcircled{1} $n=1$ \qquad On pose pour $X_1\in\mathrm{Ob}\,\unsure{\mathcal F_n|F_n}$
\[ E_1(X_1)=\partial\widetilde X_1\simeq\underset{\substack{\rotatebox{90}{$=$}\\\mathrm{arcs}(X_1)}}{\pi_0(\partial\widetilde X_1)}\deleted{\unsure{\sim|=}}\longrightarrow\pi_0(\widetilde X_1) \]
Les fibres de ce morphisme sont toutes de cardinal 2, d'où \deleted{\unsure{fonction}} \inserted{permutation} $\sigma_0$ commutant à ce morphisme, ss pt fixe (``permutation caractéristique'' de ce morphisme). Si $X_1$ est une $1$-variété, l'appl.
\diagram{\begin{tikzcd}[ampersand replacement=\&]E_1(X_1)=\partial\widetilde X_1\arrow[d]\\X_0\end{tikzcd}}
a\deleted{\unsure{st}} aussi ses fibres de card. 2, \deleted{\unsure{la}} perm. car. sera notée $\sigma_1$

\textcircled{2} Cas $n\geq2$. On pose
\[ E_n(X_n)=E_{n-1}^\circ(\partial\widetilde X_n)\qquad(=E_{n-1}(\partial\widetilde X_n)) \]
Muni des $\sigma_0,\text{---},\sigma_{n-1}$ i.e. op. de $G_{n-1}$ par hyp. de récurrence. Fonctorialité évidente, \inserted{comme composé de foncteurs $\mathcal F_n\xrightarrow{\partial\widetilde{\phantom X}}\mathcal F_{n-1}^\circ\xrightarrow{E_{n-1}^\circ}(G_{n-1}\text{-ens})$}
Si $X_n$ est une $n$-variété, \deleted{\unsure{critère}} \inserted{l'application}
\diagram{\begin{tikzcd}[ampersand replacement=\&] E_{n-1}(\partial\widetilde X_n)\arrow[d]\\E_{n-1}(X_{n-1})\end{tikzcd}}
% slanted
\margin{vérifier}
(c'est \unsure{une|un} $G_{n-2}$-hom, i.e. commute à $\sigma_1\ldots\sigma_{n-2}$) a des fibres de cardinal 2, l'hom can. \deleted{\unsure{ass.}} noté $\sigma_n$, il commute à $\sigma_1,\ldots,\sigma_{n-2}$. OK.
Fonctorialité claire encore ds $\mathcal F_n^\circ$
% archive sheet number: 37
